\section*{Asesmen}
\addcontentsline{toc}{section}{Asesmen}

Asesmen Bab VI mencakup Tugas Teori (8\%) dan bagian P1-B. Fokusnya
adalah kemampuan mengonversi dan meminimalkan otomata.

\begin{table}[H]
\centering
\caption{Instrumen asesmen Bab VI}
\label{tab:ch06-asesmen}
\begin{tabularx}{\textwidth}{@{}lXl@{}}
\toprule
\textbf{Instrumen} & \textbf{Indikator} & \textbf{Bobot} \\
\midrule
Tugas Teori & DFA/NFA, subset construction, minimalisasi & 8\% \\
P1-B & DFA token dan prototipe scanner & bagian P1 \\
Kuis & Epsilon closure dan simulasi & formatif \\
\bottomrule
\end{tabularx}
\end{table}

Kriteria kelulusan: mahasiswa dapat mengonversi NFA menjadi DFA
secara manual, meminimalkan DFA, dan menjelaskan hubungannya dengan
pengenalan token.
